\section{When splitting gains}\label{sec:theory}

\subsection{Cost model}\label{sec:costmodel}

Splitting with a selection fraction $p_0$ builds nested sets $S_l=\{\xi\ge\xi_l\}$ with $\mu(S_{l+1}\mid S_l)\approx p_0$, and one level costs $n_0(1-p_0)$ evaluations. Denote by $\pi_l=\mu(W^*\mid S_l)$ the fraction of the target in $S_l$ and by $e_l=\pi_{l+1}/\pi_l\in[0,1/p_0]$ the enrichment at level $l$; $e_l<1$ means that the coordinate is deceptive at that level.

\begin{proposition}[cost model]\label{prop:cost}
Assume that the chains sample the uniform distribution on each $S_l$ exactly, that the number of levels is unbounded, and that the final stage samples $S_K$ independently. \textup{(i)}~If the coordinate is ideal, $e_l=1/p_0$, then $K=\lceil\ln(1/p)/\ln(1/p_0)\rceil$ and $N=O(n_0\ln(1/p))$, i.e., $s=0$. \textup{(ii)}~If $e_l=p_0^{-\zeta}$ with $0<\zeta\le1$, then $K\approx\ln(1/p)/(\zeta\ln(1/p_0))$, and the cost is still logarithmic. \textup{(iii)}~If the coordinate is informative only down to a depth $K^*$ ($e_l=1$ for $l>K^*$), then $N\approx n_0+n_0(1-p_0)K^*+\ln(1/\delta)/\pi_{K^*+1}$: the exponent returns to $s=1$, with a constant-factor gain $\prod_{l\le K^*}e_l$. \textup{(iv)}~If $e_l<1$, the same count gives a factor below 1, i.e., a loss.
\end{proposition}

The proposition follows from counting levels: the target becomes frequent when $\prod_{l<K}e_l\approx1/p$, and detection in the final stage with confidence $1-\delta$ needs $\ln(1/\delta)/\pi_K$ independent samples. It is a model of the mechanism, not of the implementation. The implementation restarts after ten levels, so the logarithmic regime is available only while $p\ge p_0^{10\zeta}$. Chain states are correlated, so a coordinate with zero informative depth gives $G<1$ rather than $G=1$ (family D1, \cref{sec:scaling_p}). Finally, the exponent $s\approx0.45$ observed on family D2 over $p=10^{-2}$--$10^{-5}$ is a crossover between the logarithmic regime and the discretization limit discussed below, not an asymptotic law. The content of the proposition is a classification: an order-of-growth gain requires a coordinate that is informative on all scales, a finite informative depth yields a constant factor, and a deceptive coordinate yields a loss. The transient duration diverges at a basin boundary, where the trajectory follows the stable manifold of the edge state; if the target basin is surrounded by such a boundary and the sets $S_l$ are nested down to it, the enrichment persists on all scales.

\subsection{Product targets: the gain is geometric in the number of units}\label{sec:product}

The most favorable case for channel B is a system of weakly coupled units with hidden attractors, whose rare target is the regime in which every unit is on its hidden attractor. For uncoupled units, \cref{thm:info} and \cref{prop:cost} give matching answers.

\begin{proposition}[product targets]\label{prop:product}
Let a system on $B=B_1^k$, equipped with the product measure, consist of $k$ uncoupled copies of a unit with region $B_1$, target basin $W_1$ of measure $p_1$, and nonescaping set $Q_1\supseteq W_1$ of measure $q>p_1$. Let the composite trajectory escape as soon as any copy escapes, and let the coordinate be its escape step. Then:
\begin{enumerate}
\item[\textup{(i)}] $p=p_1^k$ and $\mu(Q_\infty)=q^k$; by \cref{cor:channels}\textup{(b)}, no algorithm that observes only the target indicator and the escape step can guarantee a gain larger than about $2.56\,q^{-k}$.
\item[\textup{(ii)}] Under the assumptions of \cref{prop:cost}, every superlevel set of the coordinate is a product set that contains the target, and the enrichment is ideal, $e_l=1/p_0$, until the superlevel set reaches $Q_\infty$, where the coordinate becomes constant.
\item[\textup{(iii)}] Splitting reaches $Q_\infty$ after $K_k=\lceil k\ln(1/q)/\ln(1/p_0)\rceil$ levels and costs
\begin{equation}
N_k\approx n_0+n_0(1-p_0)K_k+\ln(1/\delta)\,(q/p_1)^k,
\label{eq:Nk}
\end{equation}
so that
\begin{equation*}
G_k=\frac{N^0}{N_k}=\frac{q^{-k}}{1+\varepsilon_k},\qquad \varepsilon_k=\frac{n_0+n_0(1-p_0)K_k}{\ln(1/\delta)}\Bigl(\frac{p_1}{q}\Bigr)^k\to0\quad(k\to\infty):
\end{equation*}
ideal splitting attains the bound of \textup{(i)} up to the factor $2.56$.
\item[\textup{(iv)}] Along the family $p=p_1^k$, the cost grows as $N_k\propto p^{-s}$, where $s=1-\ln(1/q)/\ln(1/p_1)$.
\end{enumerate}
\end{proposition}

\begin{proof}
(i) The basin of the composite target is $W_1^k$, and the composite trajectory stays bounded if and only if every copy does, so $Q_\infty=Q_1^k$. (ii) The composite escape step is the minimum of the escape steps of the copies, so $\{\text{escape step}\ge t\}=\prod_{v=1}^k\{\text{escape step of copy } v\ge t\}$. Every such set contains $Q_\infty\supseteq W_1^k$, so $\pi_l=p/\mu(S_l)$ and $e_l=\mu(S_l)/\mu(S_{l+1})=1/p_0$ for exact quantile levels. On $Q_\infty$, the escape step is constant, and no level separates its points. (iii) The levels end when $p_0^{K}\le q^k$; the final stage then samples $Q_\infty$, where the target fraction is $(p_1/q)^k$, and needs $\ln(1/\delta)(q/p_1)^k$ samples for confidence $1-\delta$. With $N^0\approx\ln(1/\delta)p_1^{-k}$, this gives \eqref{eq:Nk} and $G_k$; $\varepsilon_k\to0$ because $K_k$ grows linearly and $(p_1/q)^k$ decays geometrically. (iv) This follows from $\ln N_k\approx k\ln(q/p_1)$ and $\ln(1/p)=k\ln(1/p_1)$.
\end{proof}

The value of escape-time information for product targets is therefore $\Theta(q^{-k})$: it grows geometrically in the number of units, with a base $1/q$ set by the nonescaping fraction of a single unit, which is measurable from a uniform sample of one unit before any search. Uniform sampling pays $p_1^{-1}$ per unit, and ideal splitting only $q/p_1$. Two remarks delimit the statement. The proposition assumes uncoupled units; for the weakly coupled copies of \cref{sec:coupled}, the measured $p$ is within 30\% of $p_1^k$, and the gains with and without coupling agree within their intervals. For units with bounded trajectories, $Q_\infty=B$, and escape times carry no information; the gains then come from the settling part $D$ of \eqref{eq:xi}, which is dominated by the units that have not yet settled, and they lie outside the scope of \cref{prop:product}.

\subsection{Diagnostics and the discretization limit}\label{sec:diagnostics}

The enrichment $e_l$ is not observable before the target is detected, but on a benchmark system with a known target, it can be estimated from a uniform sample. The decisive quantity is spatial convergence toward the target rather than the target fraction in the upper set $S_\alpha$, formed by the fraction $\alpha$ of points with the largest $\xi$. Let $d(\alpha)$ be the median distance from the nontarget points of $S_\alpha$ to the nearest target point of the sample, relative to its value for the whole sample ($\alpha=1$). Splitting gains if $d(\alpha)$ decreases as $\alpha$ decreases, because the chains move by a local random walk and need a connected path of ever-longer transients toward the basin. If $d(\alpha)$ does not decrease, long trajectories are scattered over the region, and there is no gain even if the target fraction in $S_\alpha$ grows (section~\ref{app:diagnostics}).

The coordinate is built from a trajectory with step $\Delta t$, so escape times shorter than one step are indistinguishable. In family D2 (\cref{sec:synth}), the escape time for $|w_0|\gg\rho$ is $\rho^2/(2\kappa|w_0|^2)$, and the coordinate is informative only for $|w_0|\lesssim\rho/\sqrt{2\kappa\Delta t}$; the size of this zone relative to the region limits the gain at small $p$ (\cref{sec:scaling_p}, section~\ref{app:smallp}). The time step is thus a parameter of the search, not only of the integration.
